% Replacement for §3.4 of two_problems_progress_report.tex, with the
% averaging step that justifies "controls all first-order feasible directions"
% made explicit.

\documentclass[11pt]{amsart}

\usepackage{amsmath,amssymb,amsthm,mathtools}
\usepackage[margin=1.1in]{geometry}
\usepackage{enumitem}
\usepackage{hyperref}
\usepackage{xcolor}

\hypersetup{
  colorlinks=true,
  linkcolor=blue!50!black,
  citecolor=blue!50!black,
  urlcolor=blue!50!black
}

\newtheorem{theorem}{Theorem}[section]
\newtheorem{proposition}[theorem]{Proposition}
\newtheorem{lemma}[theorem]{Lemma}
\newtheorem{conjecture}[theorem]{Conjecture}
\newtheorem{corollary}[theorem]{Corollary}

\theoremstyle{definition}
\newtheorem{definition}[theorem]{Definition}
\newtheorem{remark}[theorem]{Remark}

\newcommand{\graphon}{\mathcal W}
\newcommand{\E}{\mathbb E}
\newcommand{\1}{\mathbf 1}
\newcommand{\R}{\mathbb R}
\newcommand{\dd}{\,d}
\newcommand{\Hom}{\operatorname{Hom}}

\title[First-order stability of $T_k$ for $k\ge 6$]%
{First-order stability of the balanced $k$-partite graphon\\
under an edge-density constraint, for $k\ge 6$:\\
the averaging step made explicit}
\author{Companion note}
\date{\today}

\begin{document}
\maketitle

\begin{abstract}
This note replaces §3.4 of \emph{Progress on two graphon optimisation problems}
with a fully spelled-out proof.  The statement is that, for the six-vertex graph
$H$ of Problem~1 and every integer $k\ge 6$, the balanced complete $k$-partite
graphon $T_k$ is first-order stable under the edge-density constraint
$t(K_2,W)=1-1/k$: no first-order feasible perturbation of $T_k$ decreases
$t(H,W)$.  The original proof reduced the question to the family
$\{S_{k,a}\}$ "by the symmetry of $T_k$".  We make the averaging step that
justifies this reduction completely explicit.
\end{abstract}

\section{Setup: admissible first-order perturbations at $T_k$}

Let $H$ be the six-vertex Problem-1 graph,
\[
   E(H) = \{13,14,23,24,35,36,45,46,56\},
\]
and let $T_k$ be the balanced complete $k$-partite graphon.  Concretely, with
the partition $[0,1] = I_1 \sqcup \cdots \sqcup I_k$ into $k$ equal intervals
$I_j = [(j-1)/k, j/k]$, we have $T_k(x,y) = 1$ iff $x,y$ lie in different parts,
and $T_k(x,y) = 0$ iff $x,y$ lie in the same part.

For $f\in L^\infty([0,1]^2)$ symmetric, write $W_t := T_k + t\, f$.  The
\emph{first-order feasibility cone at $T_k$ under the edge-density constraint}
$t(K_2,W) = 1-1/k$ is the closed convex cone
\begin{equation}
\label{eq:feas-cone}
\mathcal F_k := \left\{
   f : [0,1]^2 \to \R\ \text{measurable, symmetric, } \|f\|_\infty<\infty
   \ \middle|\
   \begin{aligned}
   &f(x,y)\ge 0\ \text{a.e.\ on within-part cells},\\
   &f(x,y)\le 0\ \text{a.e.\ on cross-part cells},\\
   &\int_{[0,1]^2} f \dd x\dd y = 0
   \end{aligned}
\right\}.
\end{equation}
The first two conditions ensure that for $t>0$ small enough, $W_t$ takes values
in $[0,1]$ (so $W_t$ is a graphon), and the third ensures that
$t(K_2,W_t)=1-1/k+0\cdot t + O(t^2)$.

\begin{remark}
Strictly, every "feasible perturbation" should be defined modulo
the equivalence of graphons under measure-preserving transformations.  All
expressions below depend only on the equivalence class, so we may work with
representatives.
\end{remark}

\section{The first variation kernel and its symmetry}

For a finite graph $H$ and a graphon $W$, the standard expansion
\[
   t(H,W) = \int_{[0,1]^{V(H)}}\prod_{ij\in E(H)} W(x_i,x_j)\prod_{i\in V(H)} \dd x_i
\]
gives the first variation at $W=T_k$ in the direction $f$:
\begin{equation}
\label{eq:first-variation}
   \delta t(H,T_k)[f]
   = \sum_{ij\in E(H)}
     \int_{[0,1]^2} \mathcal K_{ij}(x,y)\, f(x,y) \dd x\dd y,
\end{equation}
where
\[
   \mathcal K_{ij}(x,y)
   = \int_{[0,1]^{V(H)\setminus\{i,j\}}}
       \prod_{\substack{i'j'\in E(H) \\ \{i',j'\}\ne\{i,j\}}} T_k(x_{i'},x_{j'})
       \prod_{v\in V(H)\setminus\{i,j\}} \dd x_v,
\]
with $x_i:=x$ and $x_j:=y$.  Symmetrising over $E(H)$, define the total kernel
\begin{equation}
\label{eq:total-kernel}
   \mathcal K(x,y) := \frac{|\mathrm{Aut}(H)|}{|\mathrm{Aut}(H)|}\cdot
   \sum_{ij\in E(H)} \mathcal K_{ij}(x,y)
   \quad (\text{symmetric in }x,y).
\end{equation}
Then
\[
   \delta t(H,T_k)[f] = \int_{[0,1]^2} \mathcal K(x,y) f(x,y) \dd x\dd y.
\]

\begin{lemma}[Block-constancy of the first-variation kernel]
\label{lem:K-block-constant}
There exist constants $c_w, c_c\in\R$ depending only on $H$ and $k$ such that
\[
   \mathcal K(x,y) =
   \begin{cases}
      c_w, & x,y\ \text{in the same part of }T_k,\\
      c_c, & x,y\ \text{in different parts of }T_k.
   \end{cases}
   \qquad\text{a.e.\ on }[0,1]^2.
\]
\end{lemma}

\begin{proof}
The kernel $\mathcal K(x,y)$ depends on $(x,y)$ only through the product
$\prod T_k(x_{i'}, x_{j'})$ for various tuples, which is an indicator of a
proper $k$-colouring condition.  Fix $(x,y)$ in the support of $\mathcal K$ and
let $j(x), j(y)\in\{1,\ldots,k\}$ be the parts containing $x$ and $y$
respectively.  For any block permutation $\sigma\in S_k$ that fixes
$j(x)$ and $j(y)$, the measure-preserving map of $[0,1]$ induced by $\sigma$
preserves $T_k$ and hence preserves $\mathcal K$.  Combined with the $S_k$-action,
this reduces the kernel value at $(x,y)$ to a function of only the pair
$\{j(x), j(y)\}$, and indeed only to the equivalence class
"same part" or "different parts" (since any two same-part pairs are related by
a block permutation, and any two different-part pairs are related by one).
\end{proof}

We will need explicit expressions for $c_w, c_c$, but for the structural
argument below they may be left abstract.

\section{The averaging step}

The space $\mathcal F_k$ in~\eqref{eq:feas-cone} is naturally acted upon by the
group $G_k$ consisting of pairs $(\sigma, \tau)$ where $\sigma\in S_k$ is a
permutation of the parts and $\tau$ is a measure-preserving automorphism of
$[0,1]$ preserving the partition $\{I_1,\ldots,I_k\}$.  Both $T_k$ and the
edge-density constraint are $G_k$-invariant, so the action preserves $\mathcal F_k$
and preserves $\delta t(H,T_k)[\cdot]$.

\begin{definition}[Symmetrised perturbation]
For $f\in\mathcal F_k$, let
\[
   \bar f(x,y) := \frac{1}{|S_k|}\sum_{\sigma\in S_k}
   \int_{\mathrm{Aut}_\partial}\ (f^{(\sigma,\tau)})(x,y)\dd\mu(\tau),
\]
where $\mathrm{Aut}_\partial$ denotes the group of partition-preserving
measure-preserving transformations of $[0,1]$ and $\mu$ is its Haar measure
(equivalently, the product over blocks of the Haar measure on the
measure-preserving group of each $I_j$).  Equivalently, $\bar f$ is the
projection of $f$ onto the $G_k$-invariant subspace of $L^2([0,1]^2)$.
\end{definition}

\begin{lemma}[Averaging preserves feasibility and the linear functional]
\label{lem:averaging}
For every $f\in\mathcal F_k$, the symmetrised perturbation $\bar f$ also lies in
$\mathcal F_k$, and
\[
   \delta t(H, T_k)[\bar f] = \delta t(H, T_k)[f].
\]
\end{lemma}

\begin{proof}
The cone $\mathcal F_k$ is defined by pointwise inequalities and a linear
constraint, all of which are $G_k$-invariant.  Hence $G_k$ preserves $\mathcal F_k$,
and so does the averaging operation, which is a convex combination of
$G_k$-images.

The second statement uses Lemma~\ref{lem:K-block-constant}: $\mathcal K$ is
$G_k$-invariant.  Therefore $\int \mathcal K\, f\dd x\dd y$ is unchanged when $f$ is
replaced by any $G_k$-image, and hence by the average $\bar f$.
\end{proof}

\begin{lemma}[Symmetric perturbations form the $S_{k,a}$ tangent line]
\label{lem:symmetric-perturbations}
Every $G_k$-invariant $\bar f\in\mathcal F_k$ has the form
\[
   \bar f(x,y) = M_w \chi_{\mathrm{within}}(x,y) - M_c \chi_{\mathrm{cross}}(x,y)
\]
for some scalars $M_w\ge 0$ and $M_c\ge 0$, with $\chi_\bullet$ the indicators
of the within-part and cross-part regions.  Moreover the edge-density
constraint $\int \bar f = 0$ forces
\[
   M_w \cdot \mathrm{vol}(\mathrm{within}) = M_c \cdot \mathrm{vol}(\mathrm{cross}),
   \quad\text{i.e.\ }\quad
   \tfrac{M_w}{k} = M_c\cdot \tfrac{k-1}{k},
   \quad
   M_c = \tfrac{M_w}{k-1}.
\]
\end{lemma}

\begin{proof}
A $G_k$-invariant symmetric function on $[0,1]^2$ must, by the same argument
as in Lemma~\ref{lem:K-block-constant}, be constant on within-part cells
and constant on cross-part cells.  The feasibility inequalities then translate
into $\bar f|_{\mathrm{within}}\ge 0$ and $\bar f|_{\mathrm{cross}}\le 0$, i.e.\
the indicated $M_w, M_c\ge 0$.  Volumes:
$\mathrm{vol}(\mathrm{within}) = k\cdot(1/k)^2 = 1/k$ and
$\mathrm{vol}(\mathrm{cross}) = 1 - 1/k$.  The constraint $\int \bar f = 0$
yields the stated relation.
\end{proof}

\begin{corollary}[Symmetric feasibility cone is one-dimensional]
\label{cor:1d-cone}
The $G_k$-invariant subset of $\mathcal F_k$ is the half-line
\[
   \{M_w\cdot \chi_{\mathrm{within}} - \tfrac{M_w}{k-1}\cdot\chi_{\mathrm{cross}}
   : M_w\ge 0\}.
\]
Up to the positive scalar $M_w$ this is the tangent direction of the family
$\{S_{k,a}\}_{a\ge 0}$ at $a=0$ (note $S_{k,0}=T_k$).
\end{corollary}

\begin{proof}
$\frac{d}{da}S_{k,a}|_{a=0}(x,y)$ equals $+1$ on within-part cells and
$-1/(k-1)$ on cross-part cells.  This is the case $M_w=1$.
\end{proof}

\section{Conclusion}

\begin{theorem}[First-order stability of $T_k$ for $k\ge 6$]
\label{thm:first-order-stability}
Let $H$ be the six-vertex Problem-1 graph.  Then for every $k\ge 6$ and every
$f\in\mathcal F_k$,
\[
   \delta t(H,T_k)[f]
   = \frac{6k^3-55k^2+142k-111}{k^5}\cdot
   \int_{\mathrm{within}} f \dd x\dd y
   \ \ge\ 0,
\]
with equality iff the symmetrisation $\bar f$ vanishes a.e.\ (equivalently,
iff $M_w(\bar f) = 0$).  In particular $T_k$ is first-order locally stable
against every edge-density-preserving feasible perturbation.
\end{theorem}

\begin{proof}
By Lemma~\ref{lem:averaging}, $\delta t(H,T_k)[f] = \delta t(H,T_k)[\bar f]$,
so we may replace $f$ by its $G_k$-symmetrisation $\bar f$.  By
Corollary~\ref{cor:1d-cone}, $\bar f = M_w\cdot \frac{d}{da}S_{k,a}|_{a=0}$
for some $M_w\ge 0$.  Hence
\[
   \delta t(H,T_k)[\bar f]
   = M_w\cdot \frac{d}{da}\,t(H,S_{k,a})\Big|_{a=0}.
\]
The progress note (§3.3, Appendix A) computes this derivative exactly:
\[
   \frac{d}{da}\,t(H,S_{k,a})\Big|_{a=0}
   = \frac{6k^3-55k^2+142k-111}{k^5}.
\]
The polynomial $6k^3-55k^2+142k-111$ evaluates to $57$ at $k=6$ and has
derivative $18k^2-110k+142$, which is positive at $k=5$ and strictly
increasing for $k\ge 5$.  Hence $6k^3-55k^2+142k-111$ is strictly increasing
for $k\ge 5$ and strictly positive for all $k\ge 6$.

Finally, $M_w\cdot \mathrm{vol}(\mathrm{within}) = M_w/k = \int_{\mathrm{within}}\bar f$,
and the symmetrisation operation is mass-preserving on the within region (more
precisely, on the convex hull of $G_k$-orbits, the within-mass of $\bar f$
equals the within-mass of $f$).  This gives the displayed identity.

Strict inequality holds whenever $M_w(\bar f)>0$, which holds iff $\bar f\not\equiv 0$ as a graphon.
\end{proof}

\begin{remark}[Why this does not settle Problem~1 for $k\ge 6$]
Theorem~\ref{thm:first-order-stability} is a \emph{negative} result for the
$k\ge 6$ part of Problem~1: it says the natural infinitesimal attack (a
perturbation of $T_k$) cannot prove that $T_k$ is not the global minimum, in
contrast with $k=3,4,5$ where the perturbation suffices.  Settling the $k\ge 6$
case therefore requires either
\begin{enumerate}[label=(\alph*)]
\item a \emph{global} construction beating $T_k$ at some $k\ge 6$ (e.g.\ a
   higher-order/non-symmetric perturbation, or a completely different
   structure such as a circular construction analogous to the $k=5$ one), or
\item a global optimality proof for $T_k$ in this range.
\end{enumerate}
A higher-order analysis is in principle the next infinitesimal step: compute
the second derivative of $t(H, S_{k,a,b})$ along a two-parameter symmetric
family (a within-block density $a$ and a single off-diagonal block density
perturbation $b$) at $(a,b)=(0,0)$.  If the resulting Hessian has a negative
eigenvalue in the constrained-feasibility cone, that already settles
non-optimality.  Otherwise we have a strict local minimum within the
two-parameter class and must broaden the search further.
\end{remark}

\begin{remark}[Comparison with $k=3,4,5$]
For $k=3,4,5$, the polynomial $6k^3-55k^2+142k-111$ evaluates to
$-26,-23,-26$ respectively, so the derivative is \emph{strictly negative},
giving an explicit edge-density-preserving feasible direction along which
$t(H,W)$ \emph{decreases} from $T_k$ at first order.  Integrating finitely
along the family $S_{k,a}$ yields the rational competitors of §3.3 of the
progress note.
\end{remark}

\end{document}
